This chapter presents the institutional context, the regulatory environment
governing anti-money laundering in Tunisia, and a review of the anomaly
detection techniques that inform the system design.

% ═══════════════════════════════════════════════════════════
\section{Host Institution}
\label{sec:amen-bank}
% ═══════════════════════════════════════════════════════════

Amen Bank is a privately held Tunisian commercial bank operating under the
supervision of the Banque Centrale de Tunisie (BCT). As of 2023, the bank
operates 150 branches distributed across 9 geographic zones (Tunis~I
through~IV, Cap~Bon, Sahel, Sfax, and South), supported by 5 Business
Centers and 11 Self-Service Areas. The institution employs 1,139 staff
members, of whom 60.5\% are assigned to the branch network and 39.5\% to
head office functions.

The internship was hosted by the \textbf{Direction Centrale du Système
d'Information} (DCSI), specifically within the \textbf{Département des
Projets Agence}, which is responsible for developing and maintaining the
information systems that support branch-level operations.

The bank's client base is composed of three economic segments: individuals
(61\% of deposits), private businesses (33\%), and institutional clients
(6\%). This composition, drawn from the audited 2023 annual report,
directly informs the client archetype design described in
Chapter~\ref{ch:synthetic}.


% ═══════════════════════════════════════════════════════════
\section{Branch Teller Operations}
\label{sec:caisse-ops}
% ═══════════════════════════════════════════════════════════

The project scope covers four teller-window (\textit{guichet}) operations
that constitute the core transactional activity of a retail bank branch:

\begin{itemize}
    \item \textbf{Retrait} (cash withdrawal) --- the client receives
          physical currency from their account. Risk exposure includes
          identity fraud, unusual withdrawal patterns, and excessive
          amounts or frequency.
    \item \textbf{Versement} (cash deposit) --- the client deposits
          physical currency into their account. This operation is the
          primary vector for structuring (\textit{smurfing}): deliberately
          splitting a large sum into multiple deposits below the reporting
          threshold.
    \item \textbf{Virement} (bank transfer) --- funds are transferred from
          the client's account to a beneficiary. Risks include transfers to
          unknown beneficiaries, duplicate transactions, and amounts
          inconsistent with the client's profile.
    \item \textbf{Remise de chèques} (cheque deposit) --- the client
          deposits one or more cheques for collection. Risks include
          cheques from unknown emitters, amounts inconsistent with
          client history, and third-party endorsement patterns.
\end{itemize}

Each operation involves two actors: the \textbf{client} initiating the
transaction and the \textbf{employee} processing it at the teller window.
Anomalies can originate from either side --- a client structuring deposits
to avoid reporting thresholds, or an employee processing an unusual volume
of high-value transactions for a concentrated set of clients. This
dual-actor nature motivates the independent client and employee profiling
described in Chapter~\ref{ch:architecture}.


% ═══════════════════════════════════════════════════════════
\section{Regulatory Framework}
\label{sec:regulatory}
% ═══════════════════════════════════════════════════════════

Tunisian banks operate under a layered anti-money laundering (LAB ---
\textit{Lutte Anti-Blanchiment}) and counter-terrorism financing (FT ---
\textit{Financement du Terrorisme}) framework involving three
institutional actors.

\subsection{Banque Centrale de Tunisie (BCT)}

The BCT issues binding circulars that define internal control obligations
for financial institutions. Circulars 2018-09, 2017-08, and 2021-05
govern LAB/FT internal controls, including the obligation to report
transactions exceeding 10,000~DT to the financial intelligence unit.
These circulars are restricted documents not publicly available; the
system's regulatory rules are therefore calibrated from the thresholds
extractable from CTAF case studies and published legislation rather than
a complete BCT circular inventory. This gap is documented explicitly
rather than papered over.

\subsection{Commission Tunisienne des Analyses Financières (CTAF)}

The CTAF is Tunisia's financial intelligence unit, responsible for
receiving and analyzing suspicious transaction reports from financial
institutions. Its 2021 annual report provides the behavioral ground truth
for what suspicious operations look like in the Tunisian context: espèces
are involved in 68\% of flagged cases, virements in 56\%, and chèques in
34\%. The concentration of reported amounts around the 10,000~DT
threshold --- with 85\% of flagged operations falling in the
5,000--20,000~DT band --- directly informed the near-threshold detection
logic in the system's decision service.

\subsection{Law n°41-2024}

Effective February 2025, this law caps individual cheques at 30,000~DT.
BCT data from Q1~2026 confirms a 24.9\% collapse in cheque usage volume
following the law's implementation. This regulatory change required
adjusting the cheque deposit archetype parameters and introduced a new
anomaly scenario: cheque amounts clustering just below the 30,000~DT
ceiling, analogous to the structuring pattern observed below the 10,000~DT
reporting threshold.


% ═══════════════════════════════════════════════════════════
\section{Existing Detection Capabilities}
\label{sec:existing-aml}
% ═══════════════════════════════════════════════════════════

The Amen Bank 2023 annual report confirms that the bank already operates
an AI-based anti-money laundering system. Based on the report's
description, this system performs \textit{bank-wide compliance reporting on
transaction streams}: it serves a compliance officer audience, operates at
a quarterly reporting tempo, and produces output aligned with
CTAF regulatory requirements.

The system developed in this internship occupies a different position in
the same regulatory umbrella. It performs \textit{real-time decision
support for branch supervisors on teller operations}: it monitors four
specific caisse operations, profiles both the client and the employee as
independent actors, operates at near-real-time latency, and uses
sequence-aware LSTM modeling to detect behavioral drift over time. The
two systems are complementary rather than competing --- the existing
system answers ``should the bank file a regulatory report?'' while this
project answers ``should the branch supervisor investigate this
transaction now?''


% ═══════════════════════════════════════════════════════════
\section{State of the Art}
\label{sec:state-of-art}
% ═══════════════════════════════════════════════════════════

This section reviews the techniques that inform the system's design,
organized from traditional approaches to deep learning methods and
explainability.

\subsection{Rule-Based Detection}

The earliest and still most widely deployed approach to transaction
monitoring relies on deterministic rules: fixed thresholds, velocity
checks, and pattern matching defined by compliance officers. Bolton and
Hand~\cite{bolton2002} provide a comprehensive survey of statistical
fraud detection methods, noting that rule-based systems dominate
production deployments due to their interpretability and auditability.

Rule-based systems have two structural limitations. First, they require
manual maintenance --- every new fraud typology demands a new rule, and
threshold calibration is labor-intensive. Second, they cannot detect
novel patterns that fall outside the predefined rule set. Despite these
limitations, regulatory requirements often mandate deterministic checks
(such as the 10,000~DT reporting threshold) that cannot be delegated to
a probabilistic model. The system developed in this project retains
deterministic rules in the Decision Service while delegating pattern
detection to the deep learning models --- combining the auditability
of rules with the adaptability of learned representations.

\subsection{Statistical Anomaly Detection}

Statistical methods model the distribution of normal behavior and flag
observations that deviate significantly. Two methods are commonly used
as baselines in the anomaly detection literature:

\textbf{Isolation Forest}~\cite{liu2008} constructs an ensemble of
random trees that partition the feature space. Anomalies, being few and
different, require fewer partitions to isolate and thus have shorter
average path lengths. The method is efficient on high-dimensional data
and does not require distributional assumptions, but it treats each
observation independently --- it cannot model temporal dependencies
between successive transactions.

\textbf{Local Outlier Factor} (LOF)~\cite{breunig2000} computes the
local density of each observation relative to its neighbors. Points in
low-density regions relative to their neighborhood are flagged as
outliers. LOF is effective for detecting local anomalies in
heterogeneous datasets but shares the same limitation as Isolation
Forest: it operates on individual data points without temporal context.

Both methods serve as useful baselines but are fundamentally
point-wise detectors. Detecting behavioral drift --- a client whose
transaction pattern gradually shifts over weeks --- requires models
that consume sequences rather than individual events.

\subsection{Autoencoders for Anomaly Detection}

An Autoencoder is a neural network trained to reconstruct its input
through a compressed latent representation. When trained exclusively
on normal data, the network learns to encode and decode the patterns
present in legitimate transactions. Anomalous inputs, which differ
from the training distribution, produce high reconstruction errors
that serve as anomaly scores.

Sakurada and Yairi~\cite{sakurada2014} demonstrated that Autoencoders
outperform linear PCA-based methods for anomaly detection in
multivariate time series, particularly when the normal data occupies a
nonlinear manifold. An and Cho~\cite{an2015} formalized the use of
reconstruction error as an anomaly score and showed that variational
Autoencoders provide calibrated uncertainty estimates alongside the
score.

In the financial domain, Autoencoders are well-suited to the class
imbalance problem: anomalies represent less than 1\% of transactions,
making supervised classification impractical without extensive labeled
data. By training only on normal transactions, the Autoencoder avoids
the need for labeled anomalies entirely --- a critical advantage in
banking environments where confirmed fraud labels are scarce and
often delayed by months.

The system uses an Autoencoder as the \textit{point anomaly detector}:
each transaction is scored independently based on how well it can be
reconstructed from the learned normal manifold.

\subsection{LSTM Networks for Sequential Anomaly Detection}

Long Short-Term Memory (LSTM) networks, introduced by Hochreiter and
Schmidhuber~\cite{hochreiter1997}, address the vanishing gradient
problem that prevents standard recurrent networks from learning
long-range dependencies. The LSTM cell's gating mechanism (input,
forget, and output gates) allows the network to selectively retain
or discard information across arbitrarily long sequences.

For anomaly detection, the LSTM is trained as a \textit{next-step
predictor}: given a sequence of recent transactions, it predicts the
features of the next transaction. The deviation between the predicted
and observed features serves as the anomaly score. Bontemps et
al.~\cite{bontemps2016} applied this approach to network intrusion
detection, showing that LSTM-based prediction models detect subtle
sequential anomalies that point-wise methods miss.

In the banking context, the LSTM captures behavioral patterns that
unfold over time: a client who normally transacts once per month but
suddenly executes five transactions in three days, or a gradual
increase in transfer amounts to a new beneficiary over several weeks.
These sequential patterns are invisible to the Autoencoder, which
evaluates each transaction in isolation.

The system uses the LSTM as the \textit{sequential anomaly detector}:
it consumes a sliding window of recent transactions per client and
flags deviations from the predicted behavioral trajectory.

\subsection{Score Fusion}

When multiple anomaly detectors operate in parallel, their scores must
be combined into a single decision. Aggarwal~\cite{aggarwal2017}
surveys fusion strategies ranging from simple averaging to learned
combination functions. The key challenge is that different detectors
may have different score distributions, making naive averaging
misleading.

The system uses a sigmoid-weighted fusion mechanism that adapts to
account maturity. New accounts with limited transaction history favor
the Autoencoder (which needs only a single event) over the LSTM
(which needs a populated sequence buffer). As the account matures
and the sequence buffer fills, the LSTM's weight increases. This
addresses the cold-start problem inherent in sequential models without
requiring a separate cold-start detection pathway.

\subsection{Explainability with SHAP}

Deploying a deep learning model in a regulated environment requires
that every alert be accompanied by an explanation that a non-technical
supervisor can understand and act upon. Lundberg and
Lee introduced SHAP (SHapley Additive
exPlanations), which attributes a model's prediction to individual
input features using game-theoretic Shapley values. SHAP provides
three properties that make it suitable for this context: local
accuracy (the attributions sum to the model output), missingness
(features not present contribute zero), and consistency (increasing a
feature's contribution never decreases its attributed importance).

For the Autoencoder, the system uses KernelSHAP --- a model-agnostic
approximation that treats the reconstruction error function as a black
box. DeepSHAP is inapplicable because the anomaly score subtracts the
input from the reconstruction, breaking the standard forward-pass
tracing that DeepSHAP requires. For the LSTM, SHAP values over the
full sequence-length-by-feature-count tensor produce attributions too
granular for a supervisor to interpret. The system instead uses an
expected-versus-actual comparison: for each interpretable feature, it
reports the LSTM's predicted value alongside the observed value,
making deviations immediately legible.

\subsection{Streaming Architectures for Real-Time Detection}

Detecting anomalies in real time requires an architecture that
processes events as they arrive rather than in periodic batches.
Kreps et al.~\cite{kreps2011} introduced the log-based streaming
paradigm with Apache Kafka, where events are written to an
append-only log partitioned by key and consumed by independent
services at their own pace. This decouples producers from consumers
and enables horizontal scaling of individual processing stages.

The system adopts this paradigm using RedPanda, a Kafka-compatible
broker that eliminates the JVM and ZooKeeper dependencies of Apache
Kafka while maintaining API compatibility. Four topics form a
processing chain: raw events are enriched with behavioral features,
scored by the ML models, evaluated by the decision service, and
written to a persistent store for dashboard consumption. Partitioning
by client identifier ensures that all transactions for a given client
are processed in order --- a requirement for the LSTM's sequential
modeling.


% ═══════════════════════════════════════════════════════════
\section{Problem Statement and Project Objectives}
\label{sec:problem-statement}
% ═══════════════════════════════════════════════════════════

The detection of anomalous teller operations at Amen Bank currently
relies on manual rules and periodic compliance reviews. This approach
cannot scale to the transaction volume of 150 branches, cannot adapt to
evolving fraud typologies without manual rule updates, and provides no
real-time decision support to branch supervisors.

The objective of this internship is to design and implement a
\textbf{real-time anomaly detection platform} that:

\begin{enumerate}
    \item Ingests teller transactions as they occur and processes them
          through a streaming pipeline with near-real-time latency.
    \item Profiles both the client and the employee as independent
          behavioral actors, detecting anomalies originating from
          either side.
    \item Combines a point anomaly detector (Autoencoder) with a
          sequential anomaly detector (LSTM) to capture both individual
          transaction deviations and behavioral drift over time.
    \item Accompanies every alert with a human-readable explanation
          generated by SHAP, enabling supervisors to make informed
          decisions without machine learning expertise.
    \item Provides a supervisor dashboard for alert review, operational
          monitoring, and transaction simulation.
    \item Implements the MLOps infrastructure --- experiment tracking,
          automated retraining, drift detection, and observability ---
          required to maintain the system beyond initial deployment.
\end{enumerate}